\documentclass[11pt,a4paper]{article}
\usepackage[T1]{fontenc}
\usepackage[margin=26mm]{geometry}
\usepackage{amsmath,amssymb,amsthm,mathtools}
\usepackage{enumitem,longtable,booktabs}
\usepackage[colorlinks=true,linkcolor=blue,citecolor=blue,urlcolor=blue]{hyperref}
\usepackage{microtype}
\setlength{\parskip}{0.3em}
\setlist{nosep}
\numberwithin{equation}{section}
\newtheorem{theorem}{Theorem}[section]
\newtheorem{lemma}[theorem]{Lemma}
\newtheorem{proposition}[theorem]{Proposition}
\newtheorem{corollary}[theorem]{Corollary}
\theoremstyle{definition}
\newtheorem{definition}[theorem]{Definition}
\theoremstyle{remark}
\newtheorem{remark}[theorem]{Remark}
\newcommand{\R}{\mathbb R}
\newcommand{\N}{\mathbb N}
\newcommand{\HH}{\mathcal H}
\newcommand{\dd}{\,\mathrm d}
\newcommand{\E}{\mathcal E}
\newcommand{\Th}{\Theta}
\newcommand{\dist}{\operatorname{dist}}
\newcommand{\dimH}{\operatorname{dim}_{\HH}}
\newcommand{\loc}{\mathrm{loc}}
\title{Compactness, tangent maps, and stratification of minimizing harmonic maps\\
\large Source reconstruction and formalization obligations for
Giaquinta--Martinazzi, printed pages 250--264}
\author{}
\date{}
\begin{document}
\maketitle

\begin{abstract}
The proof architecture is Luckhaus interpolation, annular comparison,
strong compactness of minimizers, joint upper semicontinuity of density,
homogeneous tangent maps, their translation-invariant spines, and a
Hausdorff covering argument for the singular strata. Definitions and proved
interfaces for ambient Sobolev maps, compact embedded targets, ball-local
minimality, normalized energy, monotonicity, affine transport and smooth
small-energy regularity are to be imported from the existing
\texttt{harmonicMaps} package in the existing small-energy regularity project.
This manuscript records mathematical statements and proof obligations; it
does not assert that the new results have been checked in Lean. Literal
source discrepancies and proposed repairs are identified explicitly.
\end{abstract}
\tableofcontents

\section{Scope, conventions, and reusable inputs}\label{sec:scope}
The source is Giaquinta--Martinazzi, second edition, Sections 10.3.2--10.3.3,
printed pp.~250--264, corresponding to PDF pages 252--266 in the supplied
file. The preceding statements of Theorems 10.15 and 10.18 and Corollary
10.17 are required context. The beginning of Section 10.4 on p.~264,
including Theorem 10.39, concerns weakly harmonic maps and is outside the
minimizer stratification argument. The sharpness example in Remark 10.19
and the nonuniqueness example in Remark 10.28 are contextual, not proof
prerequisites.

Fix an open set $\Omega\subset\R^n$, $n\ge2$, and a nonempty compact smooth
embedded submanifold $N\subset\R^p$, without boundary, with the induced
metric. Energy uses the Hilbert--Schmidt norm and no factor $1/2$.
All balls are open; $\overline{B_r(a)}\subset\Omega$ is the admissibility
condition. Maps and their displayed distributional gradients are identified
almost everywhere. A regular point admits a smooth representative on a
neighborhood. Statements about dilation and translation of a Sobolev map
mean equality of almost-everywhere classes, rather than an arbitrary
pointwise choice at its singular points.

\subsection{Sobolev data, local minimality, and energy}\label{def:foundation}
Use the existing $W^{1,2}_{\loc}$ predicate with its explicit weak gradient.
The existing zero-boundary predicate is the closure of smooth compactly
supported maps in the function-and-gradient $L^2$ norm. Local minimality is
\begin{equation}\label{eq:minimality}
 \E(u;B_r(a))\le\E(v;B_r(a)),\qquad
 v-u\in W^{1,2}_0(B_r(a),\R^p),\quad v\in W^{1,2}(B_r(a),N),
\end{equation}
where $\E(u;A)=\int_A|Du|^2\dd x$. This is ball-local minimality;
no boundary trace interface is implicit in the definition. A trace/gluing
theorem must connect annular boundary equalities to this zero-boundary
predicate before applying minimality.

\subsection{Monotonicity and density}\label{def:density}
Write $\Th_u(a,r)=r^{2-n}\E(u;B_r(a))$. The imported monotonicity formula is
\begin{equation}\label{eq:monotonicity}
 \Th_u(a,R)-\Th_u(a,r)
 =2\int_{B_R(a)\setminus B_r(a)}
 |x-a|^{2-n}|\partial_{\mathrm{rad},a}u|^2\dd x\ge0.
\end{equation}
In particular the density is the finite nonnegative limit
\begin{equation}\label{eq:density}
 \Th_u(a)=\lim_{r\downarrow0}\Th_u(a,r)
          =\inf_{0<r<R_a}\Th_u(a,r),
 \qquad \overline{B_{R_a}(a)}\subset\Omega.
\end{equation}
Independence of $R_a$, finiteness and the limit characterization must be
proved before using density. The upstream energy is extended nonnegative
real-valued; conversion to real energy requires its local finiteness theorem.

\subsection{Regularity and the preliminary singular-set estimate}
\label{lem:regular-density}
Imported smooth small-energy regularity supplies a uniform
$\varepsilon_*>0$, depending only on $n,N$. Smoothness implies density zero.
The printed Theorem 10.15 itself gives H\"older regularity; the smooth
conclusion used here is the already established imported bootstrap interface.
Conversely density zero yields a ball of energy below $\varepsilon_*$, hence
a smooth representative near its center. Thus
\begin{equation}\label{eq:density-gap}
 a\notin\Sigma(u)\iff\Th_u(a)=0,\qquad
 a\in\Sigma(u)\Longrightarrow\Th_u(a)\ge\varepsilon_*.
\end{equation}
The set of regular points is open, so $\Sigma(u)$ is relatively closed.
The uniform threshold, rather than a point-dependent threshold, is needed
for the final three-dimensional isolation conclusion.

\begin{lemma}[Source Corollary 10.17; preliminary estimate]
\label{lem:singular-null}
For a locally minimizing map, $\HH^{n-2}(\Sigma(u))=0$.
\end{lemma}
\begin{proof}[Proof obligations]
Equation~\eqref{eq:density-gap} places $\Sigma(u)$ in the positive
$(n-2)$-density set of the locally integrable function $|Du|^2$.
The measure-theoretic result cited as Proposition 9.21 in the book gives
zero $\HH^{n-2}$ measure of that set. Its covering/differentiation proof
must be imported from Mathlib if available, or proved. Smooth small-energy
regularity alone is not this Hausdorff-null theorem.
\end{proof}

\section{Sphere Sobolev maps and Luckhaus interpolation}
\subsection{Source Definition 10.20}\label{def:sphere-sobolev}
For $v\in L^2(S^{n-1},\R^p)$ define its degree-zero extension
$\widetilde v(r\omega)=v(\omega)$ for $r>0$. The source defines
$W^{1,2}(S^{n-1},\R^p)$ by requiring $\widetilde v\in W^{1,2}(U,\R^p)$
on some neighborhood $U$ of $S^{n-1}$. Target-valuedness is almost
everywhere. For a cylinder $S^{n-1}\times[a,b]$, use the radial extension
in the sphere variable and Sobolev regularity in both variables. The closed
interval notation includes traces of a Sobolev map on its interior.
An equivalent chart-based construction requires an equivalence theorem.
The sphere gradient $\nabla_S$, product gradient $\nabla$ and endpoint
traces require actual definitions and access lemmas.

\subsection{Source Lemma 10.21 and equations (10.39)--(10.41)}
\label{lem:luckhaus}
Put
\[
 A=\int_{S^{n-1}}(|\nabla_Su|^2+|\nabla_Sv|^2)\dd\HH^{n-1},\qquad
 D=\int_{S^{n-1}}|u-v|^2\dd\HH^{n-1}.
\]
The source asserts a constant $C$ and an ambient-valued connector $w$ with
traces $u,v$ and the joint estimates
\begin{align}
 \int_{S^{n-1}\times[0,\varepsilon]}|\nabla w|^2\dd\HH^n
 &\le C\varepsilon A+C\varepsilon^{-1}D,\label{eq:luckhaus-energy}\\
 \dist(w,N)^2
 &\le C\varepsilon^{1-n}A^{1/2}D^{1/2}
       +C\varepsilon^{-n}D\quad\text{a.e.}\label{eq:luckhaus-distance}
\end{align}
Both estimates must concern the same $w$ and one uniform $C$.
The distance estimate is essential for subsequent tubular projection.
The connector need not itself be $N$-valued.

\paragraph{Source discrepancy.}
The printed quantifier is every $\varepsilon>0$. For an arbitrary compact
target this range is false: let $N=\{0,1\}\subset\R$ and let $u=0,v=1$.
Then $A=0$, while almost every vertical Sobolev slice must connect $0$ to
$1$ and crosses the middle of the gap. The right side of
\eqref{eq:luckhaus-distance} tends to zero as $\varepsilon\to\infty$.
The construction on p.~252 uses $0<\varepsilon<1/8$, and all downstream
uses require small $\varepsilon$. A restricted small-cylinder statement is
a proposed repair; no exact Lean statement is approved here. Uniformity of
$C$ in $u,v$ and its geometric dependencies also needs an explicit statement
audit; the printed lemma does not specify these dependencies.

\subsection{The circle construction}\label{step:luckhaus-circle}
For $n=2$, choose absolutely continuous representatives on $S^1$. The
one-dimensional estimate is
\begin{equation}\label{eq:circle}
 \sup_{S^1}|u-v|^2\le C A^{1/2}D^{1/2}+CD.
\end{equation}
It follows by comparing a function to its average, the fundamental theorem
of calculus and Cauchy--Schwarz. Define
$w(\omega,s)=(1-s/\varepsilon)u(\omega)+(s/\varepsilon)v(\omega)$.
Its tangential and vertical derivatives give
\eqref{eq:luckhaus-energy}; its distance to the endpoint $u(\omega)\in N$
is at most $|u(\omega)-v(\omega)|$. For $\varepsilon\le1$,
\eqref{eq:circle} implies \eqref{eq:luckhaus-distance}.

\subsection{The higher-dimensional construction and missing estimates}
\label{step:luckhaus-skeleton}
For $n\ge3$, the source uses radial extensions to a punctured cube,
absolutely continuous representatives on almost every coordinate line,
and a translated grid of cubes of side $\varepsilon$. Interpolate on the
one-skeleton. On each $k$-face times $[0,\varepsilon]$, $k\ge2$, extend
the already specified boundary data radially from the face-cylinder center.
The extension takes place in dimension $k+1\ge3$, where its degree-zero
gradient is square integrable near the center. Choose a cube boundary by
Fubini and transfer it to the sphere by a bi-Lipschitz radial map.

\paragraph{Coordinate slicing and grid selection.}\label{lem:coordinate-acl}
Restrictions of ambient Sobolev maps to almost every coordinate line have
absolutely continuous representatives with the corresponding weak partial
derivative. Fubini also selects admissible translates and face restrictions.
These coordinate-slicing facts and the quantitative grid choice must be
proved or imported for the displayed Sobolev carrier; radial polar
integration alone does not supply them.

\paragraph{Cube-boundary/sphere transport.}\label{lem:cube-sphere}
For the source-selected $r\in[1/8,1/4]$, on the cube boundary
$x\mapsto x/|x|$ has inverse
$\omega\mapsto r\omega/\|\omega\|_\infty$ for the cube $[-r,r]^n$.
The positive lower bound on the denominators makes these maps
bi-Lipschitz. Chartwise Sobolev pullback, surface change of variables
and cylinder endpoint-trace compatibility transfer both gradient
energies and pointwise distance estimates with dimension-dependent
constants uniform in that fixed range of $r$. The geometric Lipschitz
bounds and the actual weak
composition/measure identities are separate obligations; neither follows
just from the polar integration interface.

\paragraph{Unproved source bridge.}\label{gap:luckhaus-estimates}
The book explicitly omits verification of \eqref{eq:luckhaus-energy} and
\eqref{eq:luckhaus-distance} in $n\ge3$, referring to Luckhaus [71] and
Simon [98]. Complete formalization requires translated-grid selection,
trace compatibility, quantitative radial-extension bounds, and transport
of both estimates to the sphere. These are new obligations, not extra
hypotheses of the compactness theorem. The cited complete proofs have not
been reconstructed in this manuscript.

\section{Good radii and annular gluing}
\subsection{Source Remarks 10.22--10.23, equations (10.42)--(10.44)}
\label{lem:good-radius}
For nonnegative integrable $g$, radial integration gives
\[
 \int_{B_b(y)\setminus B_a(y)}g\dd x
 =\int_a^b\int_{\partial B_s(y)}g\dd\HH^{n-1}\dd s.
\]
Chebyshev's inequality controls the exceptional radii. For $a=\rho/2$,
$b=\rho$, the exceptional set where the surface integral exceeds
$2(\theta\rho)^{-1}\int_{B_\rho\setminus B_{\rho/2}}g$ has length at
most $\theta\rho/2$. Apply this to energies and $|u-v|^2$, and intersect
their good sets. For almost every $s$, $u_s(\omega)=u(y+s\omega)$ is
sphere-Sobolev and
\begin{equation}\label{eq:spherical-scaling}
 \int_{S^{n-1}}|\nabla_Su_s|^2\dd\HH^{n-1}
 \le s^{3-n}\int_{\partial B_s(y)}|Du|^2\dd\HH^{n-1}.
\end{equation}
The exceptional-set constants in the source's (10.44) need separate care
in $n=2$: the maximum of $s^{3-n}$ on this interval changes endpoint.
The general interval estimate with a dimension-dependent constant suffices.

\subsection{Annular comparison required by compactness}
\label{lem:annular-gluing}
Let $A_{\rho,\varepsilon}(y)=B_{(1+\varepsilon)\rho}(y)
\setminus\overline{B_\rho(y)}$. The comparison needed downstream is:
for each $\Lambda>0$ there are $\delta_0>0,C<\infty$, depending only
on $n,N,\Lambda$, such that for every $y\in\R^n$, $\rho>0$,
$u,v\in W^{1,2}(A_{\rho,\varepsilon}(y),N)$ and
$0<\varepsilon\le\delta_0$, if
\[
 \rho^{2-n}\int_A(|Du|^2+|Dv|^2)\dd x\le\Lambda,\qquad
 \varepsilon^{-2n}\rho^{-n}\int_A|u-v|^2\dd x\le\delta_0,
\]
there is an $N$-valued Sobolev connector $w$ with inner trace $u$,
outer trace $v$, and
\begin{equation}\label{eq:annular-energy}
 \rho^{2-n}\int_A|Dw|^2\dd x
 \le C\rho^{2-n}\int_A(|Du|^2+|Dv|^2)\dd x
 +C\varepsilon^{-2}\rho^{-n}\int_A|u-v|^2\dd x.
\end{equation}
This is a proposed corrected comparison lemma, not a verbatim declaration
of Corollary 10.24: the printed statement assigns $u$ to both boundaries.

\begin{proof}[Construction and obligations; source (10.45)--(10.48)]
Take $\delta_0\le1/8$. Choose a common good
$s\in(\rho,(1+\varepsilon/4)\rho)$ so that
the sphere energy is at most $C\Lambda/\varepsilon$ and the sphere
$L^2$ discrepancy is at most
$C\varepsilon^{-1}\rho^{-n}\int_A|u-v|^2$.
Apply the joint Luckhaus estimates with cylinder thickness
$\varepsilon/4$. They place the connector within the tubular neighborhood:
its squared distance is bounded by
$C\sqrt{\Lambda\delta_0}+C\delta_0$ under the displayed hypotheses.
Choose $\delta_0$ small relative to the imported tubular radius. Compose
with the imported projection $\Pi$, using its derivative bound.
Transfer the cylinder to the middle annulus
$s<|x-y|<(1+\varepsilon/4)s$ using cylinder coordinate $|x-y|/s-1$.
Preserve $u$ for $\rho<|x-y|\le s$. In the next strip
$(1+\varepsilon/4)s<|x-y|<(1+\varepsilon/2)s$, use
$v(y+\psi(|x-y|)(x-y)/|x-y|)$ with
$\psi(t)=2t-(1+\varepsilon/2)s$. Its endpoint values are $s$ and
$(1+\varepsilon/2)s$, and $|\psi'|=2$.
Since $(1+\varepsilon/2)s<(1+\varepsilon)\rho$, continue with the
original $v$ to the outer boundary. The ratios $\psi(t)/t$ and the radial
Jacobian are bounded above and away from zero independently of $\rho$.
Sobolev composition and matching-trace gluing prove admissibility and
\eqref{eq:annular-energy}.

The printed (10.46) uses $\delta_0^2$ although the hypothesis uses
$\delta_0$. A square-root smallness estimate suffices, as displayed above.
The printed condition $t|\psi'(t)|\le2$ is not scale invariant for a
length-valued radial map. A uniform bound on $|\psi'|$ and on $\psi(t)/t$
is the required obligation. Neither discrepancy is an approved correction.
\end{proof}

\subsection{Traces, extension, and the imported minimizer interface}
\label{lem:trace-gluing}
The sphere connector and the annular competitor must have genuine endpoint
traces. Matching traces must imply that piecewise definitions are Sobolev
with their piecewise gradients. A map equal to $u$ in an outer collar must
differ from $u$ by an element of the imported $W^{1,2}_0$ space on the ball.
Conversely a competitor with zero-boundary difference on a smaller ball
must extend by $u$ to a larger ball. These assertions require trace,
extension and density lemmas; simply writing pointwise boundary values
does not provide them.

\section{Strong compactness of minimizing maps}
\subsection{Source Theorem 10.25}\label{thm:compactness}
Let $u_k\in W^{1,2}(\Omega,N)$ be minimizing harmonic maps and assume
\begin{equation}\label{eq:equi-energy}
 \forall\overline{B_R(a)}\subset\Omega,\qquad
 \sup_k\E(u_k;B_R(a))<\infty.
\end{equation}
There is a strictly increasing subsequence index $\kappa:\N\to\N$ and
a locally minimizing $u\in W^{1,2}_{\loc}(\Omega,N)$ such that
$u_{\kappa(k)}\to u$ strongly in $W^{1,2}_{\loc}(\Omega,\R^p)$.
The source writes global $W^{1,2}$ for its inputs; the local-input version
needed for blow-ups is an extension with the same local proof, and must be
identified as such when selecting exact declarations.

\subsection{Weak extraction and one diagonal subsequence}
\label{step:weak-extraction}
Compactness of $N$ bounds the maps in $L^2$ on bounded balls.
Rellich compactness and weak sequential compactness of the $L^2$ gradients
yield strong local map convergence and weak local gradient convergence.
Use a countable relatively compact exhaustion and nested subsequences,
then diagonalize once. Distributional derivative identities pass to the
limit; a further almost-everywhere subsequence and closedness of $N$ give
$u\in N$ almost everywhere. A separate subsequence for each uncountable
ball is not the theorem. Extension from exhaustion pieces to all compact
subdomains is a required localization lemma.

\subsection{Small-energy shell selection}\label{step:small-shell}
Fix $\overline{B_t(y)}\subset\Omega$, $0<s<t$, and a tolerance $\eta>0$.
Partition $B_t(y)\setminus\overline{B_s(y)}$ into sufficiently many
disjoint concentric shells. Uniform energy bounds and weak lower
semicontinuity bound the total energy of $u_k$ and $u$. Pigeonhole gives
one shell $A=B_{(1+\varepsilon)r}(y)\setminus\overline{B_r(y)}$,
$s\le r<(1+\varepsilon)r\le t$, with both energies at most $\eta$
along an infinite subsubsequence. Choose the partition fine enough that
$\varepsilon$ meets the annular lemma. The normalized bound $\Lambda$
and its associated $C,\delta_0$ are fixed independently of $k$ and
$\eta$, using $r\ge s$ and the total energy bound. For this fixed shell,
strong $L^2$ convergence meets its discrepancy hypothesis for large $k$.
This realizes the role of source (10.49)--(10.51) without relying on the
printed overlapping-shell indices or the ambiguous bound
$\varepsilon<(1-\varepsilon)/M$.

\subsection{Minimality passes to the limit}\label{step:limit-minimality}
Let $v$ be an arbitrary admissible competitor for $u$ on $B_s(y)$.
Extend it by $u$ to $B_r(y)$ and join it to $u_k$ through the shell.
Outside the shell use $u_k$. The trace/zero-boundary lemma makes this an
admissible competitor for $u_k$. Minimality and annular comparison imply
\begin{equation}\label{eq:comparison}
 \E(u_k;B_r(y))\le\E(v;B_s(y))
 +\E(u;B_r(y)\setminus B_s(y))+C\eta+o_k(1).
\end{equation}
Weak lower semicontinuity, cancellation of the unchanged annulus and
$\eta\downarrow0$ yield $\E(u;B_s(y))\le\E(v;B_s(y))$.
The competitor, ball and tolerance are arbitrary, so $u$ is locally
minimizing. The shell and subsubsequence may depend on the competitor;
the conclusion concerns the already fixed weak limit $u$.

\subsection{Energy convergence and strong gradient convergence}
\label{step:strong-convergence}
With $v=u$, \eqref{eq:comparison} bounds energy on $B_s$ by energy of
$u$ on $B_t$, plus $C\eta+o_k(1)$. To obtain a full limsup, apply the
shell argument to every subsequence: a subsequence violating this bound
would have a subsubsequence satisfying the comparison and contradict it.
Thus
\[
 \limsup_k\E(u_k;B_s(y))\le\E(u;B_t(y)).
\]
Let $t\downarrow s$, using absolute continuity of the energy integral and
zero volume of the boundary sphere. Combine with weak lower semicontinuity
to get equality of energy limits. In the $L^2$ Hilbert space,
\begin{equation}\label{eq:hilbert-strong}
 \|Du_k-Du\|_2^2=\|Du_k\|_2^2+\|Du\|_2^2
              -2\langle Du_k,Du\rangle\longrightarrow0.
\end{equation}
Together with strong map convergence this proves the stated local
compactness. A liminf energy estimate alone, as printed in (10.54), does
not justify this conclusion; the subsequence/limsup argument is essential.

\paragraph{Assembly of Theorem 10.25.}\label{proof:compactness-assembly}
The extraction in Section~\ref{step:weak-extraction} supplies one subsequence
and one Sobolev limit. Section~\ref{step:limit-minimality} proves this
limit is locally minimizing, and Section~\ref{step:strong-convergence}
proves strong convergence of the maps and their actual weak gradients
on every relatively compact region. These conclusions concern the same
subsequence and limit and together prove Theorem 10.25.

\section{Joint upper semicontinuity and persistence of singular points}
\subsection{Source Proposition 10.26}\label{prop:density-usc}
Suppose locally minimizing $u_k$ have locally equibounded energies,
$u_k\rightharpoonup u$ locally in $L^2$, and $y_k\to y\in\Omega$.
Then
\begin{equation}\label{eq:usc}
 \limsup_k\Th_{u_k}(y_k)\le\Th_u(y).
\end{equation}
The source writes weak $L^2$ without specifying locality; the local
interpretation is sufficient for the displayed ball proof. Compactness
and uniqueness of the prescribed weak limit give strong local $W^{1,2}$
convergence of the entire sequence: every subsequence has a strongly
convergent subsubsequence and its limit must be $u$.
For $r,h>0$ small enough and large $k$,
$B_r(y_k)\subset B_{r+h}(y)$, hence
\[
 \Th_{u_k}(y_k)\le r^{2-n}\E(u_k;B_{r+h}(y)).
\]
Pass to the strong energy limit, let $h\downarrow0$, then $r\downarrow0$.
Both the moving-center inclusion and the order of these limits are used.

\subsection{Singular points survive strong convergence}
\label{lem:singular-persistence}
If $y_k\in\Sigma(u_k)$, the uniform density gap and
\eqref{eq:usc} give $\Th_u(y)\ge\varepsilon_*>0$, so $y\in\Sigma(u)$.
For sets with $\Th_{u_k}(y_k)>\alpha>0$, the surviving estimate is
$\Th_u(y)\ge\alpha$, not necessarily $>\alpha$. Positivity is enough
for the later isolation contradiction.

\section{Blow-ups and tangent maps}
\subsection{Rescaling, source (10.56)}\label{def:blowup}
For $y\in\Omega$, $\rho>0$, define
$\Omega_{y,\rho}=\{x:y+\rho x\in\Omega\}$ and
$u_{y,\rho}(x)=u(y+\rho x)$. Its actual weak gradient is
$Du_{y,\rho}(x)=\rho Du(y+\rho x)$. The transport obligations are
\begin{equation}\label{eq:rescaling}
 \Th_{u_{y,\rho}}(x,r)=\Th_u(y+\rho x,\rho r),\qquad
 \Th_{u_{y,\rho}}(x)=\Th_u(y+\rho x).
\end{equation}
Restriction and transport of every admissible ball competitor prove local
minimality on $\Omega_{y,\rho}$. The upstream unit-ball transport theorem
must be localized to arbitrary balls before it supplies this claim.

\paragraph{Affine transport of regularity.}\label{lem:affine-regularity}
For an invertible dilation or translation $F$, smooth representatives of
$u$ pull back to smooth representatives of $u\circ F$, and conversely
through $F^{-1}$. Null sets are preserved in both directions.
Thus regular and singular loci transform by $F^{-1}$. For a map invariant
almost everywhere under a dilation or translation, the same locus is
invariant. This connects Sobolev affine identities to the singular-ray
and singular-spine arguments, without a pointwise representative at a
singular point.

\subsection{Source Proposition 10.27; existence on expanding domains}
\label{prop:tangent-existence}
For every prescribed sequence $\rho_k\downarrow0$, one can extract a
subsequence with a limit $\varphi\in W^{1,2}_{\loc}(\R^n,N)$ that is
locally minimizing. The source explicitly states existence of one such
sequence; the prescribed-sequence form is the bridge used by Lemma 10.36.
On each fixed $B_R(0)$, monotonicity bounds
\[
 \E(u_{y,\rho_k};B_R(0))=R^{n-2}\Th_u(y,\rho_kR)
 \le R^{n-2}\Th_u(y,r_0)
\]
once $\rho_kR\le r_0$. Apply local compactness on an expanding exhaustion,
diagonalize, and identify the limits on overlaps. No values outside
$\Omega_{y,\rho_k}$ are assigned by a fallback definition.
Strong energy convergence gives, for every $R>0$,
\begin{equation}\label{eq:tangent-density}
 R^{2-n}\E(\varphi;B_R(0))=\Th_u(y)=\Th_\varphi(0).
\end{equation}
Monotonicity then forces the radial derivative to vanish. A Sobolev
radial-slice lemma gives
$\varphi(\lambda\,\cdot)=\varphi$ almost everywhere for every $\lambda>0$.
This is degree-zero homogeneity. A tangent map is any limit obtained by
this process; it is not a distinguished choice and is not assumed unique.

\subsection{Source Remark 10.29}\label{prop:tangent-regularity}
The following are equivalent: $u$ is regular at $y$; $\Th_u(y)=0$;
there exists a constant tangent map at $y$. Use
\eqref{eq:density-gap} and \eqref{eq:tangent-density}. Zero energy on
every ball forces a tangent map to be almost everywhere constant on the
connected domain $\R^n$. Consequently every tangent map at a singular
point is nonconstant.

\paragraph{Zero-gradient rigidity.}\label{lem:zero-gradient}
On a ball, the imported Sobolev Poincar\'e inequality shows that zero
gradient energy makes a map almost everywhere equal to its ball average.
On overlapping balls these constants agree because the intersection has
positive volume. A chain of overlapping balls in the connected domain
$\R^n$ gives one global constant. This is a separate lemma to prove from
the Poincar\'e interface; it is not part of the definition of a Sobolev map.

\section{The spine of a homogeneous minimizer}
\subsection{Source Proposition 10.30 and (10.58)--(10.62)}
\label{prop:spine}
Let $\varphi$ be a locally minimizing degree-zero homogeneous map on
$\R^n$. Define the density-equality set
\begin{equation}\label{def:spine}
 S(\varphi)=\{y:\Th_\varphi(y)=\Th_\varphi(0)\}.
\end{equation}
Then $\Th_\varphi(y)\le\Th_\varphi(0)$ for all $y$; $S(\varphi)$ is
a real linear subspace; and for every $y\in S(\varphi)$,
$\varphi(\,\cdot+y)=\varphi$ almost everywhere.

\begin{proof}
Let the inner radius tend to zero in monotonicity about $y$. By monotone
convergence of its nonnegative integral,
\[
 \Th_\varphi(y)+2\int_{B_R(y)}|x-y|^{2-n}
       |\partial_{\mathrm{rad},y}\varphi|^2\dd x
 =R^{2-n}\E(\varphi;B_R(y)).
\]
Homogeneity about $0$ and $B_R(y)\subset B_{R+|y|}(0)$ give
\[
 R^{2-n}\E(\varphi;B_R(y))
 \le(1+|y|/R)^{n-2}\Th_\varphi(0).
\]
Let $R\to\infty$. This gives the density inequality. Equality of the
densities makes the radial derivative about $y$ vanish everywhere almost
everywhere. Thus $\varphi$ is homogeneous also about $y$:
\[
 \varphi(y+\lambda(x-y))=\varphi(x)\quad\text{a.e. in }x,
 \qquad\lambda>0.
\]
Combining the two center homogeneities gives
$\varphi(x)=\varphi(x+ty)$ with $t=\lambda-\lambda^{-1}$;
these $t$ range over all real numbers. Translation invariance preserves
local energy and therefore the density at its translated center. Sums and
scalar multiples of invariant directions are invariant, so
$t_1y_1+t_2y_2\in S(\varphi)$ whenever $y_1,y_2\in S(\varphi)$.
Also $0\in S(\varphi)$. These facts construct the genuine subspace.
\end{proof}
The printed p.~260 formula $\varphi(y+\lambda x)=\varphi(y)$ cannot
hold for a nonconstant tangent map. The centered homogeneity formula in
the proof above is the required mathematical bridge, not a literal copy.
Affine maps preserve null sets; this must be used when composing
almost-everywhere equalities.

\subsection{Source Remark 10.31}\label{lem:spine-singular}
If $\varphi$ is nonconstant, then $0\in\Sigma(\varphi)$ by tangent
regularity applied to the self-similar map. Translation invariance implies
\begin{equation}\label{eq:spine-singular}
 S(\varphi)\subset\Sigma(\varphi).
\end{equation}
Using an almost-everywhere class avoids inferring discontinuity of an
arbitrarily chosen representative from nonconstancy of that representative.

\section{Strata and their maximal index}
\subsection{Source Definition 10.32}\label{def:strata}
For integers $j\ge0$ define
\[
 S_j(u)=\{y\in\Sigma(u):
       \dim S(\varphi)\le j\text{ for every tangent map at }y\}.
\]
The universal quantifier over tangent maps is essential; replacing it by
existence of one suitable tangent would change the strata.
The spine dimension is the finite real vector-space dimension. The
Hausdorff dimension of a subset of Euclidean space is a different notion.

\subsection{Source Lemma 10.33}\label{lem:strata-top}
For $n\ge3$,
$S_0(u)\subset\cdots\subset S_{n-3}(u)
=S_{n-2}(u)=S_{n-1}(u)=\Sigma(u)$.
Indeed every tangent at a singular point is nonconstant. Its spine lies in
its singular set, whose $\HH^{n-2}$ measure is zero by
Lemma~\ref{lem:singular-null}. A subspace of dimension at least $n-2$
has positive $\HH^{n-2}$ measure on a bounded disk, a contradiction.
Thus every such spine has dimension at most $n-3$.
This is the only codimension input in the stratification argument; it
must not be replaced by the desired $n-3$ singular dimension theorem.

\section{Concentration near a tangent spine}
\subsection{Source Lemma 10.36}\label{lem:concentration}
For $0\le j\le n$, $y\in S_j(u)$ and $\delta>0$ there exists
$\varepsilon=\varepsilon(u,y,\delta)>0$ such that, for every
$0<\rho\le\varepsilon$ admissible at $y$, some $j$-dimensional linear
subspace $L_{y,\rho}$ satisfies
\begin{equation}\label{eq:concentration}
 \rho^{-1}\bigl(\{x\in B_\rho(y):
   \Th_u(x)\ge\Th_u(y)-\varepsilon\}-y\bigr)
 \subset\{z:\dist(z,L_{y,\rho})<\delta\}.
\end{equation}
The same $\varepsilon$ controls the radius and density tolerance; $L$
may depend on $\rho$. The conclusion is not uniform in $y$.

\begin{proof}
Negation supplies radii $\rho_k\downarrow0$ and tolerances
$\varepsilon_k\downarrow0$ for which every $j$-plane fails. The
prescribed-sequence blow-up bridge gives a tangent $\varphi$. Choose a
$j$-plane $L_0$ containing $S(\varphi)$, possible because its dimension
is at most $j$. Work on the compact set
$K=\overline{B_1(0)}\cap\{z:\dist(z,L_0)\ge\delta\}$ inside a larger
ball where convergence holds. If $K$ is empty the inclusion is immediate.
Otherwise upper semicontinuity and the density-equality definition give
a strict uniform gap: for some $a>0$,
$\Th_\varphi(z)\le\Th_\varphi(0)-2a$ on $K$.
Failure for $L_0$ yields $z_k\in K$ with
$\Th_{u_{y,\rho_k}}(z_k)\ge\Th_u(y)-\varepsilon_k$.
Extract $z_k\to z\in K$. Joint upper semicontinuity and
\eqref{eq:tangent-density} force
$\Th_\varphi(z)\ge\Th_\varphi(0)$, a contradiction.
\end{proof}
The source already uses $\overline{B_1(0)}$ in (10.65) and for the limit
point $x_0$. Its strict density bound away from the spine must also hold
on this closed ball; the global density maximum and the density-equality
definition supply that fact. The slightly smaller tolerance $a$ avoids
losing a strict density gap in the limit. Finite-dimensional extension of
the spine to a $j$-plane is an independent linear-algebra obligation.

\paragraph{Extension of the spine.}\label{step:spine-plane}
For any subspace $V\subset\R^n$ with $\dim V\le j\le n$, extend a basis
of $V$ to a basis of $\R^n$ and take the span of its first $j$ vectors,
retaining all the basis vectors of $V$. This constructs a $j$-plane
containing $V$. Apply this lemma to the characterized spine subspace;
the plane is an existence witness, not a canonical choice.

\section{Flat approximation and Hausdorff coverings}
\subsection{Source Definition 10.37}\label{def:flatness}
A set $A\subset\R^n$ has the $\delta$-approximation property at dimension
$j$ if there exists $r_0>0$ such that, for every $a\in A$ and
$0<r\le r_0$, some $j$-plane $L_{a,r}$ satisfies
\[
 A\cap B_r(a)\subset\{x:\dist(x,a+L_{a,r})<\delta r\}.
\]
The plane after undoing rescaling is affine and passes through $a$.
No rectifiability or quantitative-stratification conclusion is part of
this source property.

\subsection{Source Lemma 10.38}\label{lem:flat-cover}
For fixed $n,j$, $0\le j\le n$, there is a positive function
$\beta(\delta)\to0$ as $\delta\downarrow0$ such that every set with
the $\delta$-approximation property satisfies
$\HH^{j+\beta(\delta)}(A)=0$.

\begin{proof}[Covering argument with explicit recentering]
For $0<\delta<1/8$, put $\sigma=4\delta<1/2$.
A $j$-dimensional disk of radius $R$ has a cover by at most
$C_j\delta^{-j}$ balls of radius $\delta R$ within its plane; $j=0$
uses one ball.
\label{step:disk-net}
For $j>0$, choose a maximal $(\delta R/2)$-separated finite set in the
closed disk. Balls of radius $\delta R/4$ about its centers are disjoint
and lie in the disk of radius $(1+\delta/4)R$. Comparing $j$-dimensional
volumes bounds their number by $(1+4/\delta)^j\le5^j\delta^{-j}$.
This packing bound ensures the greedy separated-set construction
terminates, and maximality gives the required disk cover. Thus one can
take $C_j=5^j$. Orthogonal projection of $A\cap B_R(a)$ onto the
approximating plane puts it in that disk. The tube is therefore covered
by ambient balls of radius $2\delta R$ with centers in the plane.
Discard balls not meeting $A\cap B_R(a)$. For each remaining ball
choose a point of $A\cap B_R(a)$ in it. Recenter there and double the
radius to $4\delta R=\sigma R$. The children now have centers in $A$,
so the approximation property can be applied at every iteration.

Write $c\ge\max(1,4^jC_j)$, so the number of children satisfies
$Q\le c\sigma^{-j}$. Set
\begin{equation}\label{eq:beta}
 \beta(\delta)=\frac{\log(2c)}{\log(1/(4\delta))}>0.
\end{equation}
Then $Q\sigma^{j+\beta(\delta)}\le1/2$ and $\beta(\delta)\to0$.
Cover a bounded part of $A$ initially by finitely many $A$-centered
balls of radius $R\le r_0$. The $q$th generation has radius $\sigma^qR$
and total $(j+\beta)$-power of the radii at most $2^{-q}T_0$.
Diameters contribute the fixed factor $2^{j+\beta}$; the maximal
diameter tends to zero. The definition of Hausdorff measure gives zero.
Unbounded sets are countable unions of bounded pieces. For
$\delta\ge1/8$, use $\beta(\delta)=n-j+1$, since $\HH^{n+1}(A)=0$.
\end{proof}
The source iteration does not explicitly justify recentering its new
balls at points of $A$. The argument above records the required bridge;
the chosen formula for $\beta$ is a sufficient construction, not an
author-fixed normalization of that existential witness.

\section{Dimension of the strata and isolated singularities}
\subsection{Source Proposition 10.34: literal statement and stronger goal}
\label{prop:stratum-dimension}
The printed statement on p.~261 gives $\dimH S_j(u)\le n-3$ for
$0\le j\le n-3$, whereas the proof on pp.~262--263 establishes the
stratified estimate
\begin{equation}\label{eq:stratum-dimension}
 \dimH S_j(u)\le j.
\end{equation}
The latter is the natural proposed formalization goal and implies the
literal statement. Selecting it as a source-facing frozen declaration
requires an explicit ruling; the discrepancy is not silently resolved.
Equivalently, for every $s>j$, prove $\HH^s(S_j(u))=0$.

\begin{proof}[Derivation recorded in the source proof]
Fix $\delta>0$. Let $S^{\delta}_{j,i}$ consist of points in $S_j(u)$
for which \eqref{eq:concentration} holds with
$\varepsilon=1/i$, with $B_{1/i}(y)$ admissible. Shrinking a valid
$\varepsilon$ preserves the inclusion, so
$S_j(u)=\bigcup_{i\ge1}S^{\delta}_{j,i}$.
For $q\ge1$ set
\[
 A^{\delta}_{i,q}=S^{\delta}_{j,i}\cap
  \{y:(q-1)/i<\Th_u(y)\le q/i\}.
\]
Positive density on the singular set makes these bins exhaustive.
For $x,y$ in the same bin,
$\Th_u(x)>\Th_u(y)-1/i$. Consequently each
$A^{\delta}_{i,q}$ has the $\delta$-approximation property with
$r_0=1/i$. The covering lemma and countable subadditivity give
$\HH^{j+\beta(\delta)}(S_j(u))=0$ for this fixed $\delta$.
For an arbitrary $s>j$, choose a sufficiently small $\delta$ with
$j+\beta(\delta)<s$, and use monotonicity of Hausdorff exponents.
This proves \eqref{eq:stratum-dimension}.
\end{proof}
For $j\le n-3$ this also proves the literal dimension assertion of
Proposition 10.34. Its other assertion is proved by the isolation
argument below.
The auxiliary sets depend on $\delta$; one cannot keep an individual
$A^{\delta}_{i,q}$ fixed and then assert it has every approximation
parameter. The order is: fix $\delta$, take a countable union, then let
$\delta$ tend to zero at the level of $S_j(u)$.

\subsection{The density superlevel isolation conclusion}
\label{prop:isolation}
For $n=3$ and $\alpha>0$, the set
$S_0(u)\cap\{y:\Th_u(y)>\alpha\}$ has no accumulation point in
$\Omega$, hence is discrete. Suppose distinct $x_k$ in this set tend
to $x_0\in\Omega$. Rescale by $r_k=|x_k-x_0|$ and extract a tangent
$\varphi$ at $x_0$ with
$(x_k-x_0)/r_k\to\xi\in S^2$. Upper semicontinuity gives
$\Th_\varphi(\xi)\ge\alpha>0$, so $\xi$ is singular.
Homogeneity transports singularity to the ray
$\{t\xi:t>0\}$. Its bounded nontrivial segment has positive
$\HH^1$ measure, contradicting
$\HH^1(\Sigma(\varphi))=0$ from Corollary 10.17.
Relative closedness of $\Sigma(u)$ and admissibility of $x_0$ must be
retained; boundary accumulation is not excluded.

\subsection{Source Corollary 10.35 and Theorem 10.18}
\label{cor:singular-dimension}
For $n\ge3$ a locally minimizing harmonic map satisfies
$\dimH\Sigma(u)\le n-3$. For $n=3$ its singular set has no
accumulation in $\Omega$. Combine $S_{n-3}(u)=\Sigma(u)$ with the
literal dimension assertion of Proposition 10.34. For isolation choose
$0<\alpha<\varepsilon_*$,
so every singular point belongs to the strict density superlevel set.
The book's choice $\alpha=\varepsilon_0$ needs a strict-gap justification;
choosing half the imported threshold avoids that issue.

\paragraph{The earlier proof of Theorem 10.18 on p.~258.}
Its three-dimensional contradiction uses the same compactness, density
semicontinuity, homogeneous blow-up and singular ray. Equation (10.55)
is exactly the constant tangent energy identity
\eqref{eq:tangent-density}. The higher-dimensional sentence refers back
to the dimension-reduction proof of Theorem 10.11; that earlier route is
not needed when Corollary 10.35 is assembled from the strata. It remains
an explicit external proof route, not an omitted premise of this route.

\section{Formalization boundaries and proposed declaration surfaces}
\label{sec:formalization}
The new proof obligations divide into four mathematical components:
\begin{enumerate}
 \item Sphere/cylinder Sobolev data, traces, Luckhaus estimates and
 annular comparison; followed by local extraction and strong compactness.
 \item Density, regular/singular definitions, preliminary Hausdorff nullity
 and joint upper semicontinuity; no $n-3$ estimate is assumed.
 \item Expanding-domain rescaling, tangent maps, homogeneity, spines and
 their dimensions. Tangent maps are a relation, not a chosen unique map.
 \item Density concentration, finite-dimensional flat covers, stratified
 Hausdorff estimates and the isolation corollary.
\end{enumerate}
Recommended small declaration surfaces are strong local compactness,
joint density semicontinuity, the joint tangent existence/energy/homogeneity
theorem, the joint spine/subspace/invariance theorem, and the final joint
stratification statement. The choice of exact carriers and quantifiers,
and every source repair above, must be resolved before elaborating and
approving full Lean declarations. Shared constants and witnesses in
Luckhaus and gluing must stay joint. Helper conclusions may not be added
as hypotheses of these public theorems.

\begin{thebibliography}{99}\small
\bibitem{GM} M. Giaquinta and L. Martinazzi,
\emph{An Introduction to the Regularity Theory for Elliptic Systems,
Harmonic Maps and Minimal Graphs}, second edition,
Sections 10.3.2--10.3.3, printed pp.~250--264.
\bibitem{Luckhaus} S. Luckhaus,
Partial H\"older continuity for minima of certain energies among maps
into a Riemannian manifold, \emph{Indiana Univ. Math. J.} 37 (1988),
349--367. Source reference [71]; quantitative proof not reconstructed here.
\bibitem{SimonMaps} L. Simon,
\emph{Theorems on Regularity and Singularity of Energy Minimizing Maps},
Birkh\"auser, 1996. Source reference [98].
\bibitem{SimonSing} L. Simon,
Singularities of geometric variational problems, in
\emph{Nonlinear Partial Differential Equations in Differential Geometry},
R. Hardt and M. Wolf (eds.), AMS, 1996. Source reference [99].
\bibitem{SU} R. Schoen and K. Uhlenbeck,
A regularity theory for harmonic maps, \emph{J. Differential Geom.}
17 (1982), 307--335. Source reference [95].
\end{thebibliography}
\end{document}
